\section{Introduction and Report Methodology}
\label{sec:intro}

Blood-flow simulation sits between continuum mechanics, numerical analysis,
and biomedical interpretation. A stenosis is a useful organizing example
because a vessel narrowing changes velocity, pressure, wall loading, and
downstream flow in ways that depend on the retained state and the observable
being compared. A three-dimensional (3D) Navier--Stokes CFD model, an
axisymmetric or two-dimensional (2D) reduction of that continuum description,
a reduced one-dimensional (1D) area--flow model, and a lumped
zero-dimensional (0D) outlet network therefore do not answer the same
mathematical question. Anatomical narrowing and
functional hemodynamic significance are also different observables:
pressure--flow measurements and coronary imaging studies distinguish lesion
appearance from functional assessment
\cite[p.~4401]{DerimayEtAl2021CoronaryStenosisObstruction}
\cite[Abstract; Methods]{VanDenHoogenEtAl2022CoronaryLesionCTA}.

This report develops a review framework for reading mathematical blood-flow
results across model tiers. It organizes the literature around the ingredients
that make a computed hemodynamic quantity meaningful: the continuum description,
model class, closures, boundary data, reported observable, and supporting
numerical or empirical evidence. The idealized stenosis computation developed
later is a worked application of that reading discipline; the literature review
remains broader than the example.

\subsection{Literature Review Scope}
\label{subsec:narrative-review-scope}

The review is narrative and scoping, not systematic. Sources are used for the
role they play in interpreting mathematical hemodynamics: derivation, taxonomy,
numerical evidence, clinical or experimental observables, reproducibility, or
learned-model scope. Appendix~\ref{app:code-and-ai-use} lists supplementary
source and reproducibility records; the main text uses those records only to
support bounded claims.

\begin{table}[!htb]
    \centering
    \small
    \caption{Primary source roles used in the narrative review.}
    \label{tab:intro-source-roles}
    \begin{tabularx}{\textwidth}{@{}
        >{\raggedright\arraybackslash}p{0.24\textwidth}
        >{\raggedright\arraybackslash}p{0.34\textwidth}
        >{\raggedright\arraybackslash}X@{}}
        \toprule
        \textbf{Source role} & \textbf{What it supports} & \textbf{Representative use} \\
        \midrule
        Derivation and theory
            & Governing equations, reductions, wall laws, coupling assumptions.
            & Continuum-to-1D reductions and multiscale coupling
            \cite{FormaggiaEtAl2003OneDimensionalBloodFlow,QuarteroniVeneziani2003GeometricalMultiscale}. \\
        Review and taxonomy
            & Model classes, state variables, and clinical or numerical context.
            & 0D/1D hierarchy and pressure--flow network interpretation
            \cite{ShiEtAl2011BloodFlowModelReview}. \\
        Numerical benchmark or reproducibility study
            & Scheme comparison, reproducibility comparison, and inspectable numerical setup.
            & 1D benchmark cases and open finite-volume solvers
            \cite{BoileauEtAl2015Benchmark,BenemeritoEtAl2024OpenBF}. \\
        Clinical or experimental observable
            & External meaning of pressure ratios, anatomy/function distinction,
            and measured wall or flow quantities.
            & Angiographic versus functional stenosis assessment
            \cite{ToninoEtAl2010FAMEAngioFunctional,Blanco2018FFR1D3D}. \\
        Learned-model study
            & Performance within a declared training target, geometry family, and
            output observable.
            & PINN, neural differential, and operator-learning scope limits
            \cite{YuEtAl2026PINNHemodynamicsReview}. \\
        \bottomrule
    \end{tabularx}
\end{table}

The inclusion criterion is relevance to the mathematical reading of blood-flow
simulation: what equations are solved, which physics are retained or closed,
what data are required, what observables are produced, and what evidence can
support a comparison. A source is included when it supports at least
one of those questions for the model hierarchy, closures, numerical evidence
standard, or 1D--3D observation operator. Derivation papers establish equations
and assumptions; benchmark papers establish comparative numerical behavior;
clinical and experimental studies address external relevance for specified
observables; and open-source or reproducibility materials document numerical
setup and reviewability. The review does not claim exhaustive coverage of all
clinical or computational hemodynamics.

\subsection{Research Problem and Questions}
\label{subsec:research-problem-questions}

The research problem is to synthesize how blood-flow models are formulated,
closed, observed, and supported by numerical evidence. The report asks three
review-oriented questions:
\begin{enumerate}
    \item What mathematical model hierarchies are used to simulate blood flow,
    and how do retained state, closure burden, coupling structure, and native
    observables distinguish 3D CFD, FSI, 2D or axisymmetric, 1D, 0D,
    multiscale, and learned formulations?
    \item How do rheology, wall mechanics, geometry, inlet and outlet data,
    numerical discretization, evidence standards, and observation operators
    affect the interpretation of reported hemodynamic quantities?
    \item What limitations remain across the literature when state variables,
    resolved scales, computational cost, practical identifiability, validation
    status, reproducibility, and uncertainty sources are compared?
\end{enumerate}

\subsection{Literature Review and Report Roadmap}
\label{subsec:objectives-contributions}
\label{subsec:intro-notation-guide}
\label{subsec:report-organization}

The report proceeds in two stages. Sections~\ref{sec:continuum-description}--%
\ref{sec:literature-synthesis} build the reading framework by fixing the
continuum vocabulary, comparing model tiers, making closure and observable
choices explicit, and separating numerical evidence categories that are often
blurred together in hemodynamics. Section~\ref{sec:case-study} then uses an
idealized stenosis example to show how those distinctions matter in practice:
the model, boundary approximation, discretization, verification evidence, and
1D--3D observation rule are specified before any comparison is interpreted. The
final discussion returns from that example to the review questions and states
the bounded methodological contribution directly.

The main continuum symbols are used in their standard units unless a generated
numerical table states otherwise: $\rho$ denotes density, $u_z$ axial velocity,
$\eta$ dynamic viscosity, $\nu=\eta/\rho$ kinematic viscosity, $R$ radius,
$A$ cross-sectional area, and $Q$ volumetric flow. The case study uses the
physical variables $A_{\mathrm{phys}}$ and $Q_{\mathrm{phys}}$ when comparing
with continuum quantities, but the implemented finite-volume state stores
$a=R^2$ and $q=Q_{\mathrm{phys}}/\pi$. Thus
$A_{\mathrm{phys}}=\pi a$, $Q_{\mathrm{phys}}=\pi q$, and the section-mean
velocity is $\bar u=q/a$. Appendix~\ref{app:domain-specific-notation} collects
the full symbol and unit table used by the report.

The rest of the manuscript follows that sequence: continuum definitions,
model hierarchy, closure choices, numerical evidence, synthesis, case-study
application, and integrated discussion.
